\subsection{Rechtevergabe}
\label{sec:lizenzen}

Die Werke eines Schöpfers (wie Texte, Musikstücke, Bilder, Videos usw.) sind normalerweise urheberrechtlich geschützt. Der Schöpfer kann aber entscheiden, dass er Werke anderen Menschen zur Verfügung stellt, ohne dass sie ausdrücklich um Erlaubnis fragen müssen. Dazu veröffentlicht er die Werke mit einem entsprechenden Hinweis, dass er zum Beispiel das Recht zum Kopieren, Verändern und Wiederveröffentlichen allen anderen zugesteht.

Für juristische Laien ist es allerdings schwierig, einen entsprechenden Rechtstext zu formulieren. Schließlich soll deutlich sein, was erlaubt ist und was nicht, und es soll auch kein Missbrauch mit den zur Verfügung gestellten Werken möglich sein (etwa, dass jemand behauptet, er selbst sei Schöpfer dieser Werke). Um diesem Problem zu begegnen, wurde die Organisation Creative Commons gegründet, um solche Rechtstexte (Lizenzen) zu erarbeiten. \cite{wikicc}



\hyperlink{https://www.wur.nl/en/article/What-are-Creative-Commons-licenses.htm}{FAQ- What are Creative Commons licenses?}

\hyperlink{https://creativecommons.org/licenses/}{CC Bilder }

\hyperlink{https://moqod-software.medium.com/understanding-open-source-and-free-software-licensing-c0fa600106c9}{Software tetxte}

\hyperlink{https://choosealicense.com/}{guide}